\documentclass[11pt,a4paper]{article}
\usepackage[utf8]{inputenc}
\usepackage[T1]{fontenc}
\usepackage[margin=1in]{geometry}
\begin{document}



% Title block (centered)
\thispagestyle{empty}
\begin{center}
    \vspace*{0.5cm}
    {\large \textbf{Tool Research Proposal}}\\[0.5cm]
    \hr\\[0.4cm]
    {\huge \bfseries \color{unalblue} \projecttitle}\\[0.4cm]
    \hr\\[1.2cm]

    \begin{tabular}{rl}
        \textbf{Candidate:} & \candidate \\ 
        \textbf{Supervisor:} & \supervisor \\ 
        \textbf{Department:} & \department \\ 
        \textbf{Institution:} & \institution \\ 
        \textbf{Date:} & \dateprep
    \end{tabular}
    \vspace{1.2cm}

    \begin{tcolorbox}[colback=boxbg, colframe=unalblue, title=Abstract, fonttitle=\bfseries]
        This report was generated by the Research Accelerant local agent. It synthesizes indexed documents and user-provided chat history into a concise LaTeX report suitable for PDF compilation.
    \end{tcolorbox}
\end{center}
\vfill
\clearpage

% --- PAGE 2: Table of Contents ---
\pagenumbering{arabic}
\tableofcontents
\clearpage




\begin{resultbox}[Definition]
In the context of a research proposal, a problem statement is a concise, clear, and evidence-based description of a specific \textbf{researchable} issue, a \textbf{demonstrable} gap in knowledge, or a \textbf{measurable} real-world challenge that your research is \textbf{designed to solve}.
\end{resultbox}

\vspace{1em}

\textbf{\large Structured Problem Statement}\par\vspace{0.5em}

\textbf{1. What is Known}\\
The literature on machine learning has established foundational knowledge across Across the 3 reviewed studies on machine learning, the following methodological patterns emerged:. Key studies have contributed empirical evidence and theoretical frameworks.

\vspace{0.5em}

\textbf{2. What is Missing}\\
Geographic scope: Most studies on machine learning may be concentrated in specific regions, limiting generalizability to diverse populations.

\vspace{0.5em}

\textbf{3. Affected Stakeholders}\\
Researchers, practitioners, and policymakers working with machine learning are directly affected by this gap, as are the populations ultimately served by evidence-based interventions.

\vspace{0.5em}

\textbf{4. Consequences of Inaction}\\
Without addressing this gap, decision-makers will continue to rely on incomplete or non-generalizable evidence, potentially leading to ineffective strategies and missed opportunities for improvement in machine learning.

\vspace{0.5em}

\textbf{5. How the Study Addresses the Gap}\\
The proposed study will employ rigorous methodology to directly investigate geographic scope: most studies on machine learning may be concentrated in specific regions, limiting generalizability to diverse populations., generating actionable evidence to fill this critical gap in the literature.

\vspace{1em}

\begin{resultbox}[Full Problem Statement]
Despite advances in understanding machine learning, a critical gap remains: geographic scope: most studies on machine learning may be concentrated in specific regions, limiting generalizability to diverse populations.. This is problematic because current evidence is insufficient to guide effective practice across diverse contexts. The proposed study addresses this gap by conducting rigorous, contextually grounded research that will produce generalizable findings and actionable recommendations for stakeholders in the field.
\end{resultbox}

\vspace{1em}

\textbf{\large Validation Checklist}\par\vspace{0.3em}

\begin{itemize}[leftmargin=*, label=$\square$]
    \item Is the problem \textbf{researchable}? (Can it be investigated through empirical methods?)
    \item Is the gap \textbf{demonstrable}? (Can you cite evidence that it exists?)
    \item Is the challenge \textbf{measurable}? (Can outcomes be quantified or assessed?)
    \item Are \textbf{stakeholders} clearly identified?
    \item Are \textbf{consequences of inaction} articulated?
    \item Does the statement tell reviewers \textbf{why}, \textbf{what}, and \textbf{how}?
\end{itemize}

\vfill
\begin{center}
\textit{Problem statement generated by Research Accelerant Agent on \today.}
\end{center}

\end{document}